\documentclass[12pt,letterpaper]{article}
\usepackage[utf8]{inputenc}
\usepackage[spanish]{babel}
\usepackage{amsmath}
\usepackage{amsfonts}
\usepackage{amssymb}
\usepackage[margin=2.5cm]{geometry}

\begin{document}

\begin{center}
    \Large \textbf{Evaluación de Examen 2: Unidad II} \\
    \large Física para Ciencias de la Salud (005-1713) \\
    \vspace{0.5cm}
    \textbf{Estudiante:} Yilselys Naranjo \quad \textbf{C.I.:} 33.385.584
    \vspace{0.3cm} \\
    \large \textbf{Calificación Final: 9/10 puntos}
\end{center}

\vspace{0.5cm}

\section*{Análisis de Resultados}

\subsection*{1. Cinemática (Aceleración media)}
\textbf{Procedimiento y resultado:} Extrae de manera estructurada los datos iniciales y plantea la ecuación de aceleración media. Efectúa la sustitución y la resta correctas, dividiendo $0,6 \text{ m/s}$ entre $0,15 \text{ s}$. El cálculo aritmético arroja el resultado analítico exacto de $4 \text{ m/s}^2$ y el manejo dimensional es correcto. \\
\textbf{Veredicto:} \textbf{Correcto} (2/2 pts).

\subsection*{2. Dinámica (Leyes de Newton)}
\textbf{Procedimiento y resultado:} Plantea la fórmula vectorial de la Segunda Ley de Newton ($\Sigma \vec{F} = m \cdot \vec{a}$) y realiza el despeje algebraico para hallar la aceleración. Sustituye los valores dividiendo $150 \text{ N}$ entre $25 \text{ kg}$, logrando el valor matemático perfecto de $6$. Sin embargo, comete un error dimensional al denotar el resultado final con unidades de fuerza ($6 \text{ N}$) en lugar de unidades de aceleración ($6 \text{ m/s}^2$). \\
\textbf{Veredicto:} \textbf{Parcialmente correcto} (Penalización por unidad incorrecta en el resultado final) (1,5/2 pts).

\subsection*{3. Trabajo Mecánico}
\textbf{Procedimiento y resultado:} Extrae los datos e identifica la fórmula general del trabajo, incluyendo de forma explícita el componente trigonométrico $\cos(\theta)$. Efectúa correctamente la multiplicación ($400 \text{ N} \cdot 0,8 \text{ m} \cdot \cos(0^\circ)$). El cálculo aritmético da exactamente $320 \text{ J}$. \\
\textbf{Veredicto:} \textbf{Correcto} (2/2 pts).

\subsection*{4. Energía Potencial}
\textbf{Procedimiento y resultado:} Extrae rigurosamente los datos del problema. Plantea el producto de las tres variables según la ecuación de la energía potencial gravitatoria ($E_p = m \cdot g \cdot h$) integrando las unidades dentro de la sustitución ($1,2 \text{ kg} \cdot 9,8 \text{ m/s}^2 \cdot 1,5 \text{ m}$). Obtiene el resultado verificado de $17,64 \text{ J}$. \\
\textbf{Veredicto:} \textbf{Correcto} (2/2 pts).

\subsection*{5. Potencia Mecánica}
\textbf{Procedimiento y resultado:} Extrae los datos y relaciona el trabajo efectuado y el tiempo planteando la fracción respectiva ($P = \frac{W}{T}$). El desarrollo de la división de $75 \text{ J} / 60 \text{ s}$ arroja la cifra exacta de $1,25$. No obstante, escribe la unidad final explícitamente como ``Joules'' en lugar de ``Watts'', lo cual es un error conceptual para la magnitud de la potencia. \\
\textbf{Veredicto:} \textbf{Parcialmente correcto} (Penalización por unidad incorrecta) (1,5/2 pts).

\vspace{0.5cm}
\section*{Comentarios Generales y Sugerencias}
¡Un muy buen examen! La estudiante demuestra un dominio aritmético perfecto, ya que logró obtener todos y cada uno de los números correctos para los cinco problemas planteados, mostrando además una excelente y limpia estructuración separando ``Datos'' y ``Solución''.
\begin{itemize}
    \item Se requiere repasar urgentemente el \textbf{análisis dimensional y las unidades del Sistema Internacional}. En física biomédica, confundir Newtons con $\text{m/s}^2$ (Pregunta 2) o Joules con Watts (Pregunta 5) puede llevar a malinterpretaciones severas de los fenómenos.
    \item Como un pequeñísimo apunte formal: los símbolos de las unidades estándar de kilos y segundos se escriben en minúscula ($\text{kg}$ en lugar de $\text{Kg}$, y $\text{s}$ en lugar de $\text{sg}$).
    \item ¡Sigue con este mismo entusiasmo y nivel de habilidad matemática en la próxima unidad, prestando un poco más de atención a las unidades finales!
\end{itemize}

\end{document}